\documentclass[10pt,a4paper]{article}

%double si on veut une version prof et une version élève. simple si on veut une seule version.
\newcommand{\buildmode}{simple}

\InputIfFileExists{version.tex}{}{}
% Version (définie par le script)
\providecommand{\version}{eleve}

\usepackage{feuilleinterro}

%%%à modifier à chaque fois

\renewcommand{\niveau}{1G}
\renewcommand{\chapitre}{Suites, généralités}
\renewcommand{\numchap}{1}
\renewcommand{\theme}{Suites}

\begin{document}

% =========================================================
% PAGE 1 : VERSION A
% =========================================================

\interroversion{A}

\textbf{Exercice 1 -- Calcul avec des fractions \hfill /2}

Calculer, en détaillant les étapes, et donner le résultat sous la forme d'une fraction irréductible :
\[
A = \frac{2}{3}\times\frac{3}{4}-\frac{1}{6}.
\]

\vspace{3cm}

\textbf{Exercice 2 -- Suite définie par récurrence \hfill /4}

On considère la suite \((u_n)\) définie pour tout entier naturel \(n\) par
\[
\left\{\begin{array}{rcl}
u_0 &=& 1 \\
u_{n+1} &=& 2u_n+1
\end{array}\right.
\]

\begin{enumerate}
\item Quel est le mode de définition de cette suite ? Justifier.

\vspace{2cm}

\item Calculer \(u_1\), \(u_2\) et \(u_4\).
\end{enumerate}

\vspace{3cm}

\textbf{Exercice 3 -- Sens de variation d'une suite \hfill /4}

On considère la suite \((u_n)\) définie pour tout entier naturel \(n\) par \(u_n=3n-5\).

Déterminer le sens de variation de la suite \((u_n)\) (on pourra calculer \(u_{n+1}-u_n\) en fonction de \(n\)).

\newpage

% =========================================================
% PAGE 2 : VERSION B
% =========================================================

\interroversion{B}

\textbf{Exercice 1 -- Calcul avec des fractions \hfill /2}

Calculer, en détaillant les étapes, et donner le résultat sous la forme d'une fraction irréductible :
\[
B = \frac{5}{6}\times\frac{2}{3}-\frac{1}{2}.
\]

\vspace{3cm}

\textbf{Exercice 2 -- Suite définie par récurrence \hfill /4}

On considère la suite \((u_n)\) définie pour tout entier naturel \(n\) par
\[
\left\{\begin{array}{rcl}
u_0 &=& 2 \\
u_{n+1} &=& 2u_n-1
\end{array}\right.
\]

\begin{enumerate}
\item Quel est le mode de définition de cette suite ? Justifier.

\vspace{2cm}

\item Calculer \(u_1\), \(u_2\) et \(u_4\).
\end{enumerate}

\vspace{3cm}

\textbf{Exercice 3 -- Sens de variation d'une suite \hfill /4}

On considère la suite \((u_n)\) définie pour tout entier naturel \(n\) par \(u_n=2n-7\).

Déterminer le sens de variation de la suite \((u_n)\) (on pourra calculer \(u_{n+1}-u_n\) en fonction de \(n\)).

\newpage

% =========================================================
% PAGE 3 : VERSION C
% =========================================================

\interroversion{C}

\textbf{Exercice 1 -- Calcul avec des fractions \hfill /2}

Calculer, en détaillant les étapes, et donner le résultat sous la forme d'une fraction irréductible :
\[
C = \frac{3}{5}\times\frac{5}{6}-\frac{1}{3}.
\]

\vspace{3cm}

\textbf{Exercice 2 -- Suite définie par récurrence \hfill /4}

On considère la suite \((u_n)\) définie pour tout entier naturel \(n\) par
\[
\left\{\begin{array}{rcl}
u_0 &=& 0 \\
u_{n+1} &=& 3u_n+2
\end{array}\right.
\]

\begin{enumerate}
\item Quel est le mode de définition de cette suite ? Justifier.

\vspace{2cm}

\item Calculer \(u_1\), \(u_2\) et \(u_4\).
\end{enumerate}

\vspace{3cm}

\textbf{Exercice 3 -- Sens de variation d'une suite \hfill /4}

On considère la suite \((u_n)\) définie pour tout entier naturel \(n\) par \(u_n=4n-3\).

Déterminer le sens de variation de la suite \((u_n)\) (on pourra calculer \(u_{n+1}-u_n\) en fonction de \(n\)).

\newpage

% =========================================================
% PAGE 4 : VERSION D
% =========================================================

\interroversion{D}

\textbf{Exercice 1 -- Calcul avec des fractions \hfill /2}

Calculer, en détaillant les étapes, et donner le résultat sous la forme d'une fraction irréductible :
\[
D = \frac{7}{10}\times\frac{5}{6}-\frac{1}{4}.
\]

\vspace{3cm}

\textbf{Exercice 2 -- Suite définie par récurrence \hfill /4}

On considère la suite \((u_n)\) définie pour tout entier naturel \(n\) par
\[
\left\{\begin{array}{rcl}
u_0 &=& 1 \\
u_{n+1} &=& 4u_n-2
\end{array}\right.
\]

\begin{enumerate}
\item Quel est le mode de définition de cette suite ? Justifier.

\vspace{2cm}

\item Calculer \(u_1\), \(u_2\) et \(u_4\).
\end{enumerate}

\vspace{3cm}

\textbf{Exercice 3 -- Sens de variation d'une suite \hfill /4}

On considère la suite \((u_n)\) définie pour tout entier naturel \(n\) par \(u_n=5n-6\).

Déterminer le sens de variation de la suite \((u_n)\) (on pourra calculer \(u_{n+1}-u_n\) en fonction de \(n\)).

\end{document}
